% AA-65
\chapter[Wagi najwyższe]{Reprezentacje półprostych algebr Liego i wagi najwyższe}
\label{aa:ch:lie-highest-weights}

Rozkład korzeniowy algebry Liego opisuje jej reprezentację dołączoną. Teraz rozważymy dowolną reprezentację skończenie wymiarową. Zobaczymy, że występowanie całkowitych wag oraz szczególna rola wag najwyższych wynikają już z teorii $\mathfrak{sl}_2$ z AA-60. Nie będziemy zakładać pełnej redukowalności reprezentacji ogólnej algebry półprostej.

W całym rozdziale $K$ jest algebraicznie domkniętym ciałem charakterystyki zero, $\mathfrak g$ niezerową skończenie wymiarową półprostą algebrą Liego nad $K$, $\mathfrak h$ jej maksymalną podalgebrą toralną, a $\Delta=(\alpha_1,\ldots,\alpha_r)$ ustaloną bazą układu korzeni. Niech $h_i=h_{\alpha_i}$ oraz $e_i,f_i$ będą znormalizowanymi generatorami z AA-61. W zapisie $xv$ rozumiemy działanie $\rho(x)v$ reprezentacji $\rho:\mathfrak g\to\mathfrak{gl}(V)$.

\section{Rozkład na przestrzenie wagowe}

\begin{definicja}
\label{aa:def:lie-module-weight}
Niech $\mathfrak g$ będzie niezerową skończenie wymiarową półprostą algebrą Liego nad algebraicznie domkniętym ciałem $K$ charakterystyki zero, $\mathfrak h$ jej maksymalną podalgebrą toralną, $\Delta=(\alpha_1,\ldots,\alpha_r)$ bazą układu korzeni, a $h_i=h_{\alpha_i}$ oraz $e_i,f_i$ znormalizowanymi generatorami z Twierdzenia~\ref{aa:tw:lie-simple-generators}. Niech $V$ będzie skończenie wymiarową reprezentacją $\mathfrak g$. Dla funkcjonału $\lambda\in\mathfrak h^*$ określamy przestrzeń wagową
\[
 V_\lambda=\{v\in V:tv=\lambda(t)v\text{ dla każdego }t\in\mathfrak h\}.
\]
Jeśli $V_\lambda\ne0$, funkcjonał $\lambda$ nazywamy wagą reprezentacji, a $\dim V_\lambda$ jej krotnością.
\end{definicja}

\begin{twierdzenie}
\label{aa:tw:lie-weight-decomposition}
Niech $\mathfrak g$ będzie niezerową skończenie wymiarową półprostą algebrą Liego nad algebraicznie domkniętym ciałem $K$ charakterystyki zero, $\mathfrak h$ jej maksymalną podalgebrą toralną, $\Delta=(\alpha_1,\ldots,\alpha_r)$ bazą układu korzeni, a $h_i=h_{\alpha_i}$ oraz $e_i,f_i$ znormalizowanymi generatorami z Twierdzenia~\ref{aa:tw:lie-simple-generators}. Niech $V$ będzie skończenie wymiarową reprezentacją $\mathfrak g$. Wówczas $V$ jest sumą prostą swoich przestrzeni wagowych, występuje tylko skończenie wiele wag, a dla każdej wagi $\lambda$ i każdego $i$ zachodzi $\lambda(h_i)\in\mathbb Z$ (liczby całkowite utożsamiamy z ich obrazem w $K$). Jeśli $x\in\mathfrak g_\alpha$ dla korzenia $\alpha$, to $xV_\lambda\subset V_{\lambda+\alpha}$. Każdy podmoduł $U\subset V$ jest sumą prostą przestrzeni $U\cap V_\lambda$.
\end{twierdzenie}
\begin{proof}
Trójka $e_i,f_i,h_i$ rozpina podalgebrę izomorficzną z $\mathfrak{sl}_2(K)$. Ograniczenie reprezentacji do niej jest na mocy Twierdzenia~\ref{aa:tw:sl2-complete-reducibility} sumą prostą modułów $L(m)$. Na każdym $L(m)$ operator $h_i$ jest diagonalizowalny i ma całkowite wartości własne $m,m-2,\ldots,-m$. Zatem operator $h_i$ na $V$ jest diagonalizowalny i wszystkie jego wartości własne są całkowite.

Operatory $h_i$ komutują, gdyż $\mathfrak h$ jest abelowa. Rozkładamy $V$ na przestrzenie własne $h_1$; pozostałe operatory zachowują te przestrzenie dzięki komutowaniu. Ich ograniczenia są diagonalizowalne, gdyż ich wielomiany minimalne dzielą wielomiany minimalne operatorów wyjściowych, które mają różne pierwiastki. Powtarzając rozkład dla kolejnych $h_i$, otrzymujemy sumę prostą wspólnych przestrzeni własnych. Ponieważ $h_i$ tworzą bazę $\mathfrak h$ na mocy Lematu~\ref{aa:lem:lie-simple-root-basis}, każdy układ wartości własnych wyznacza jeden funkcjonał $\lambda$. Są to dokładnie przestrzenie wagowe.

Dla $t\in\mathfrak h$, $v\in V_\lambda$ i $x\in\mathfrak g_\alpha$ mamy
\[
 t(xv)=[t,x]v+x(tv)=(\alpha(t)+\lambda(t))xv.
\]
Jeśli $\lambda+\alpha$ nie jest wagą, prawa przestrzeń wagowa jest zerowa. Wreszcie podmoduł $U$ jest niezmienniczy względem wszystkich $h_i$, więc te same rozkłady na przestrzenie własne stosują się do ograniczeń na $U$. Ich wspólne przestrzenie własne to $U\cap V_\lambda$.
\end{proof}

\begin{definicja}
\label{aa:def:lie-dominant-integral}
Niech $\mathfrak g$ będzie niezerową skończenie wymiarową półprostą algebrą Liego nad algebraicznie domkniętym ciałem $K$ charakterystyki zero, $\mathfrak h$ jej maksymalną podalgebrą toralną, $\Delta=(\alpha_1,\ldots,\alpha_r)$ bazą układu korzeni, a $h_i=h_{\alpha_i}$ oraz $e_i,f_i$ znormalizowanymi generatorami z Twierdzenia~\ref{aa:tw:lie-simple-generators}. Niech $\omega_1,\ldots,\omega_r\in\mathfrak h^*$ będą funkcjonałami określonymi przez $\omega_j(h_i)=\delta_{ij}$. Nazywamy je wagami fundamentalnymi. Krata wag i zbiór wag całkowitych dominujących to odpowiednio
\[
 P=\sum_{i=1}^r\mathbb Z\omega_i,\qquad
 P_+=\sum_{i=1}^r\mathbb Z_{\ge0}\omega_i.
\]
Dla $\lambda,\mu\in\mathfrak h^*$ piszemy $\mu\preceq\lambda$, gdy
$\lambda-\mu=\sum_i n_i\alpha_i$ z $n_i\in\mathbb Z_{\ge0}$.
\end{definicja}

Wagi fundamentalne tworzą bazę przestrzeni $\mathfrak h^*$, ponieważ są bazą sprzężoną do $(h_i)$. Relacja $\preceq$ jest porządkiem częściowym: zwrotność i przechodniość wynikają z definicji, a antysymetria z liniowej niezależności $\alpha_i$ i charakterystyki zero. Nie porządkujemy elementów samego ciała $K$. Z Twierdzenia~\ref{aa:tw:lie-weight-decomposition} każda waga reprezentacji należy do $P$.

\section{Symetria wag}

\begin{twierdzenie}
\label{aa:tw:lie-weight-reflection}
Niech $\mathfrak g$ będzie niezerową skończenie wymiarową półprostą algebrą Liego nad algebraicznie domkniętym ciałem $K$ charakterystyki zero, $\mathfrak h$ jej maksymalną podalgebrą toralną, $\Delta=(\alpha_1,\ldots,\alpha_r)$ bazą układu korzeni, a $h_i=h_{\alpha_i}$ oraz $e_i,f_i$ znormalizowanymi generatorami z Twierdzenia~\ref{aa:tw:lie-simple-generators}. Niech $V$ będzie skończenie wymiarową reprezentacją $\mathfrak g$. Dla każdego $i$ i $\lambda\in\mathfrak h^*$ zachodzi
\[
 \dim V_\lambda=\dim V_{s_i\lambda},
 \qquad s_i\lambda=\lambda-\lambda(h_i)\alpha_i.
\]
W szczególności zbiór wag wraz z krotnościami jest niezmienniczy względem grupy generowanej przez odbicia $s_i$ (grupy Weyla).
\end{twierdzenie}
\begin{proof}
Ustalmy $i$ i oznaczmy $\mathfrak h'=\ker\alpha_i\subset\mathfrak h$. Ponieważ $\alpha_i(h_i)=2$, mamy $\mathfrak h=Kh_i\oplus\mathfrak h'$. Dla $\eta=\lambda|_{\mathfrak h'}$ rozważmy przestrzeń
\[
 W_\eta=\{v\in V:tv=\eta(t)v\text{ dla każdego }t\in\mathfrak h'\}.
\]
Operatory $e_i,f_i,h_i$ zachowują $W_\eta$, gdyż dla $t\in\mathfrak h'$ komutują z działaniem $t$: $[t,e_i]=[t,f_i]=[t,h_i]=0$. Stosujemy zatem Twierdzenie~\ref{aa:tw:sl2-complete-reducibility} do reprezentacji $\mathfrak{sl}_2(K)$ na $W_\eta$. Każdy jej składnik $L(m)$ ma takie same krotności wartości własnych $a$ i $-a$ operatora $h_i$, gdyż ciąg $m,m-2,\ldots,-m$ jest symetryczny (obie krotności mogą być zerowe). Ta równość zachodzi również w sumie prostej.

Przestrzeń własna $h_i$ na $W_\eta$ dla wartości $\lambda(h_i)$ to $V_\lambda$, ponieważ $\mathfrak h=Kh_i\oplus\mathfrak h'$. Funkcjonał $s_i\lambda$ ma na $\mathfrak h'$ ograniczenie $\eta$, a na $h_i$ wartość $\lambda(h_i)-2\lambda(h_i)=-\lambda(h_i)$. Stąd teza o krotnościach. Iterowanie równości daje niezmienniczość względem wszystkich iloczynów odbić.
\end{proof}

\section{Wektory najwyższej wagi}

\begin{definicja}
\label{aa:def:lie-highest-vector}
Niech $\mathfrak g$ będzie niezerową skończenie wymiarową półprostą algebrą Liego nad algebraicznie domkniętym ciałem $K$ charakterystyki zero, $\mathfrak h$ jej maksymalną podalgebrą toralną, $\Delta=(\alpha_1,\ldots,\alpha_r)$ bazą układu korzeni, a $h_i=h_{\alpha_i}$ oraz $e_i,f_i$ znormalizowanymi generatorami z Twierdzenia~\ref{aa:tw:lie-simple-generators}. Niech $V$ będzie skończenie wymiarową reprezentacją $\mathfrak g$. Niezerowy wektor $v\in V_\lambda$ nazywamy wektorem najwyższej wagi $\lambda$, gdy $e_i v=0$ dla wszystkich $1\le i\le r$.
\end{definicja}

Warunek ten jest równoważny $\mathfrak n^+v=0$, gdzie $\mathfrak n^+=\bigoplus_{\alpha>0}\mathfrak g_\alpha$. Istotnie, $\mathfrak n^+$ jest generowana przez $e_i$ na mocy Twierdzenia~\ref{aa:tw:lie-simple-generators}, a jeśli $xv=yv=0$, to $[x,y]v=x(yv)-y(xv)=0$.

\begin{lemat}
\label{aa:lem:lie-highest-cyclic}
Niech $\mathfrak g$ będzie niezerową skończenie wymiarową półprostą algebrą Liego nad algebraicznie domkniętym ciałem $K$ charakterystyki zero, $\mathfrak h$ jej maksymalną podalgebrą toralną, $\Delta=(\alpha_1,\ldots,\alpha_r)$ bazą układu korzeni, a $h_i=h_{\alpha_i}$ oraz $e_i,f_i$ znormalizowanymi generatorami z Twierdzenia~\ref{aa:tw:lie-simple-generators}. Niech $V$ będzie skończenie wymiarową reprezentacją $\mathfrak g$, a $v$ wektorem najwyższej wagi $\lambda$. Podmoduł $W$ generowany przez $v$ jest rozpięty przez wektory
\[
 f_{i_1}\cdots f_{i_k}v,\qquad k\ge0,
\]
gdzie dla $k=0$ otrzymujemy $v$. Każda waga $\mu$ podmodułu $W$ spełnia $\mu\preceq\lambda$, a $W_\lambda=Kv$.
\end{lemat}
\begin{proof}
Niech $S$ będzie rozpiętością wskazanych wektorów. Każdy niezerowy taki wektor ma wagę $\lambda-\sum_{j=1}^k\alpha_{i_j}$, na mocy Twierdzenia~\ref{aa:tw:lie-weight-decomposition}, więc $S$ jest niezmiennicza względem $\mathfrak h$. Jest także niezmiennicza względem $f_i$ z definicji. Aby wykazać niezmienniczość względem $e_i$, indukujemy po długości słowa. Dla długości zero mamy $e_i v=0$. Dla słowa $f_jw$ korzystamy z relacji z AA-61:
\[
 e_i f_jw=f_j e_iw+\delta_{ij}h_iw.
\]
Pierwszy składnik należy do $S$ na mocy założenia indukcyjnego i niezmienniczości względem $f_j$, drugi dzięki niezmienniczości względem $h_i$. Generatory $e_i,f_i$ generują $\mathfrak g$, stąd $S$ jest podmodułem zawierającym $v$. Każdy podmoduł zawierający $v$ zawiera $S$, zatem $W=S$.

Dla słowa niepustego suma $\sum_j\alpha_{i_j}$ nie może być zerem, ponieważ $\alpha_i$ są liniowo niezależne nad $K$ i charakterystyka jest zero. Z rozkładu na przestrzenie wagowe wynika więc, że tylko słowo puste przyczynia się do składowej $W_\lambda$. Otrzymujemy $W_\lambda=Kv$ oraz tezę o porządku wag.
\end{proof}

\begin{twierdzenie}
\label{aa:tw:lie-highest-necessary}
Niech $\mathfrak g$ będzie niezerową skończenie wymiarową półprostą algebrą Liego nad algebraicznie domkniętym ciałem $K$ charakterystyki zero, $\mathfrak h$ jej maksymalną podalgebrą toralną, $\Delta=(\alpha_1,\ldots,\alpha_r)$ bazą układu korzeni, a $h_i=h_{\alpha_i}$ oraz $e_i,f_i$ znormalizowanymi generatorami z Twierdzenia~\ref{aa:tw:lie-simple-generators}. Niech $V$ będzie niezerową skończenie wymiarową reprezentacją $\mathfrak g$. Istnieje w niej wektor najwyższej wagi, a każda waga takiego wektora należy do $P_+$. Jeśli reprezentacja $V$ jest nieprzywiedlna, to istnieje dokładnie jedna waga najwyższa $\lambda$, jej przestrzeń ma wymiar jeden, a wszystkie wagi $\mu$ reprezentacji spełniają $\mu\preceq\lambda$.
\end{twierdzenie}
\begin{proof}
Skończony niepusty zbiór wag ma element maksymalny względem $\preceq$; wybieramy taki $\lambda$. Dla każdego $i$ funkcjonał $\lambda+\alpha_i$ jest od niego ściśle większy, zatem nie jest wagą. Każdy niezerowy $v\in V_\lambda$ spełnia więc $e_iv=0$.

Dla dowolnego wektora najwyższej wagi $v$ ustalmy $i$. Relacje $\mathfrak{sl}_2$ dają
\[
 h_i f_i^jv=(\lambda(h_i)-2j)f_i^jv,\qquad
 e_i f_i^jv=j(\lambda(h_i)-j+1)f_i^{j-1}v
\]
dla $j\ge1$, jak w Lemacie~\ref{aa:lem:sl2-commutation}. Niezerowe $f_i^jv$ są liniowo niezależne ze względu na różne wartości własne $h_i$. Istnieje więc pierwszy indeks $m+1$ taki, że $f_i^{m+1}v=0$. Stosując $e_i$, otrzymujemy
$(m+1)(\lambda(h_i)-m)f_i^mv=0$. Ponieważ $K$ ma charakterystykę zero i $f_i^mv\ne0$, mamy $\lambda(h_i)=m\ge0$. Dla wszystkich $i$ daje to $\lambda\in P_+$.

Jeżeli reprezentacja $V$ jest nieprzywiedlna, podmoduł generowany przez $v$ jest równy $V$. Lemat~\ref{aa:lem:lie-highest-cyclic} daje $V_\lambda=Kv$ i $\mu\preceq\lambda$ dla każdej wagi $\mu$. Jeśli $w$ jest innym wektorem najwyższej wagi $\nu$, to również generuje $V$, więc analogicznie $\lambda\preceq\nu$. Antysymetria daje $\nu=\lambda$, a $w\in Kv$.
\end{proof}

\begin{twierdzenie}[Jednoznaczność reprezentacji]
\label{aa:tw:lie-highest-uniqueness}
Niech $\mathfrak g$ będzie niezerową skończenie wymiarową półprostą algebrą Liego nad algebraicznie domkniętym ciałem $K$ charakterystyki zero, $\mathfrak h$ jej maksymalną podalgebrą toralną, $\Delta=(\alpha_1,\ldots,\alpha_r)$ bazą układu korzeni, a $h_i=h_{\alpha_i}$ oraz $e_i,f_i$ znormalizowanymi generatorami z Twierdzenia~\ref{aa:tw:lie-simple-generators}. Niech $V,W$ będą niezerowymi skończenie wymiarowymi nieprzywiedlnymi reprezentacjami $\mathfrak g$. Są izomorficzne wtedy i tylko wtedy, gdy mają tę samą wagę najwyższą.
\end{twierdzenie}
\begin{proof}
Izomorfizm zachowuje przestrzenie wagowe i warunek $e_iv=0$, więc konieczność wynika z Twierdzenia~\ref{aa:tw:lie-highest-necessary}.

Załóżmy, że obie wagi najwyższe są równe $\lambda$. Wybierzmy niezerowe wektory najwyższej wagi $v\in V$, $w\in W$. W sumie prostej $V\oplus W$ z działaniem składowym wektor $(v,w)$ ma wagę najwyższą $\lambda$. Niech $C$ będzie generowanym przez niego podmodułem. Z Lematu~\ref{aa:lem:lie-highest-cyclic} wiemy, że $C_\lambda=K(v,w)$. Projekcje $p:C\to V$, $q:C\to W$ są homomorfizmami modułów i surjekcjami: ich obrazy zawierają odpowiednio $v,w$, a moduły są nieprzywiedlne.

Gdyby $\ker p\ne0$, to $\ker p\subset 0\oplus W$ byłby niezerowym podmodułem $0\oplus W$, więc równym $0\oplus W$. Wtedy $(0,w)\in C_\lambda$, co jest sprzeczne z $C_\lambda=K(v,w)$, bo $v\ne0$. Zatem $p$ jest izomorfizmem. Identyczny argument dla $q$ daje drugi izomorfizm. Złożenie $q\circ p^{-1}$ jest żądanym izomorfizmem $V\to W$.
\end{proof}

Otrzymaliśmy konieczność dominującej całkowitości i jednoznaczność. Nie dowodziliśmy jeszcze, że każda waga z $P_+$ występuje jako waga najwyższa skończenie wymiarowej reprezentacji nieprzywiedlnej. Nie należy też utożsamiać nieprzywiedlności z nierozkładalnością bez twierdzenia o pełnej redukowalności.

% Zachowujemy przykład i odpowiadający mu diagram na jednej stronie.
\newpage
\section{Przykład: reprezentacja naturalna i sprzężona}

\begin{przyklad}
\label{aa:prz:lie-highest-sl3}
Niech $\mathfrak g=\mathfrak{sl}_3(K)$, $\mathfrak h$ będzie podalgebrą diagonalną, a $e_1=E_{12}$, $e_2=E_{23}$, $f_1=E_{21}$, $f_2=E_{32}$, $h_1=E_{11}-E_{22}$, $h_2=E_{22}-E_{33}$. Współrzędne wagi zapisujemy jako $(\lambda(h_1),\lambda(h_2))$. Na standardowej bazie $b_1,b_2,b_3$ reprezentacji naturalnej $K^3$ wagi to
\[
 (1,0),\quad (-1,1),\quad (0,-1).
\]
Wektor $b_1$ jest zerowany przez $e_1,e_2$, a $f_1b_1=b_2$, $f_2b_2=b_3$. Reprezentacja jest nieprzywiedlna: każdy niezerowy podmoduł zawiera wektor bazowy dzięki rozkładowi wagowemu i jednowymiarowości przestrzeni wagowych; stosując odpowiednie $e_i$ uzyskujemy $b_1$, a następnie $f_i$ całą bazę. Waga najwyższa to $\omega_1$.

Na przestrzeni sprzężonej $(K^3)^*$ działanie określa wzór $(x\varphi)(v)=-\varphi(xv)$. Jest to reprezentacja, gdyż komutator dwóch takich operatorów działający na $\varphi$ i obliczony w $v$ daje $-\varphi([x,y]v)$. Dla bazy sprzężonej $b_1^*,b_2^*,b_3^*$ wagi są przeciwne do poprzednich:
$(-1,0),(1,-1),(0,1)$. Wektor $b_3^*$ jest zerowany przez $e_1,e_2$, a $f_2b_3^*=-b_2^*$ i $f_1b_2^*=-b_1^*$. Podobnie jak wyżej każdy niezerowy podmoduł zawiera któryś wektor bazowy; zastosowanie $e_i$ prowadzi do $b_3^*$, a $f_i$ generują pozostałe. Reprezentacja sprzężona jest nieprzywiedlna o wadze najwyższej $\omega_2$. Obie reprezentacje mają ten sam wymiar, lecz nie są izomorficzne na mocy Twierdzenia~\ref{aa:tw:lie-highest-uniqueness}.
\end{przyklad}

\begin{center}
\begin{tikzpicture}[every node/.style={font=\small}]
\node (a) at (0,2.4) {$(1,0)$};
\node (b) at (0,1.2) {$(-1,1)$};
\node (c) at (0,0) {$(0,-1)$};
\node (d) at (5,2.4) {$(0,1)$};
\node (e) at (5,1.2) {$(1,-1)$};
\node (f) at (5,0) {$(-1,0)$};
\draw[->] (a)--node[right] {$f_1$} (b);
\draw[->] (b)--node[right] {$f_2$} (c);
\draw[->] (d)--node[right] {$f_2$} (e);
\draw[->] (e)--node[right] {$f_1$} (f);
\node at (0,-.6) {$K^3$};
\node at (5,-.6) {$(K^3)^*$};
\end{tikzpicture}

\small Wagi i niezerowe działania operatorów obniżających wagę w obu reprezentacjach.
\end{center}
